5. Encuentra una base ortogonal para \( \mathbb{R}^{4} \) que contenga a los vectores \( (2,1,0,-1) \) y \( (1,0,3,2) \). \\

%%SOLO MEJORAR ESCRITURA Y TENER CUIDADO EN EL PROCESO DE
\begin{solution}
Para encontrar una base ortogonal para \( \mathbb{R}^{4} \) que contenga a los vectores \( \mathbf{v}_1 = (2, 1, 0, -1) \) y \( \mathbf{v}_2 = (1, 0, 3, 2) \), podemos usar el proceso de ortogonalización de Gram-Schmidt. Este proceso nos permitirá transformar un conjunto de vectores linealmente independientes en un conjunto de vectores ortogonales que generan el mismo subespacio. Primero, verificamos que los vectores dados sean linealmente independientes. Para esto, calculamos el determinante de la matriz formada por estos vectores. Dado que solo tenemos dos vectores, si ninguno es un múltiplo escalar del otro, son linealmente independientes. Claramente, \( \mathbf{v}_1 \) y \( \mathbf{v}_2 \) no son múltiplos escalares, por lo que podemos proceder. \\
Ortogonalización de Gram-Schmidt. \\
1. Vector \( \mathbf{u}_1 \):
\[
   \mathbf{u}_1 = \mathbf{v}_1 = (2, 1, 0, -1)
   \]
2. Vector \( \mathbf{u}_2 \):
\[
   \mathbf{u}_2 = \mathbf{v}_2 - \frac{\mathbf{v}_2 \cdot \mathbf{u}_1}{\mathbf{u}_1 \cdot \mathbf{u}_1} \mathbf{u}_1
   \]
Calculamos los productos escalares:
\[
   \mathbf{v}_2 \cdot \mathbf{u}_1 = 1 \cdot 2 + 0 \cdot 1 + 3 \cdot 0 + 2 \cdot (-1) = 2 - 2 = 0
   \]
\[
   \mathbf{u}_1 \cdot \mathbf{u}_1 = 2^2 + 1^2 + 0^2 + (-1)^2 = 4 + 1 + 0 + 1 = 6
   \]
\[
   \therefore \mathbf{u}_2 = \mathbf{v}_2 - 0 \cdot \mathbf{u}_1 = \mathbf{v}_2 = (1, 0, 3, 2)
   \]

Hasta aquí, tenemos que los vectores \( \mathbf{u}_1 \) y \( \mathbf{u}_2 \) son ortogonales.
Paso 3: Completar la base ortogonal
Ahora necesitamos encontrar dos vectores adicionales para completar una base de $\mathbb{R^4}$. Podemos elegir vectores canónicos, por ejemplo, y ortogonalizarlos respecto a los vectores obtenidos.


Vector \( \mathbf{v}_3 = (0, 1, 0, 0) \):
\[
   \mathbf{u}_3 = \mathbf{v}_3 - \frac{\mathbf{v}_3 \cdot \mathbf{u}_1}{\mathbf{u}_1 \cdot \mathbf{u}_1} \mathbf{u}_1 - \frac{\mathbf{v}_3 \cdot \mathbf{u}_2}{\mathbf{u}_2 \cdot \mathbf{u}_2} \mathbf{u}_2
   \]
Calculamos:
\[
   \mathbf{v}_3 \cdot \mathbf{u}_1 = 0 \cdot 2 + 1 \cdot 1 + 0 \cdot 0 + 0 \cdot (-1) = 1
   \]
\[
   \mathbf{v}_3 \cdot \mathbf{u}_2 = 0 \cdot 1 + 1 \cdot 0 + 0 \cdot 3 + 0 \cdot 2 = 0
   \]
Entonces:
\[
   \mathbf{u}_3 = \mathbf{v}_3 - \frac{1}{6} \mathbf{u}_1 = \left(0, 1, 0, 0\right) - \frac{1}{6} \left(2, 1, 0, -1\right) = \left(-\frac{1}{3}, \frac{5}{6}, 0, \frac{1}{6}\right)
   \]
Calculamos el producto punto de $\textbf{v}_4$ con $\textbf{u}_1$​:
\[
   \mathbf{v}_4 \cdot \mathbf{u}_1 = 0 \cdot 2 + 0 \cdot 1 + 1 \cdot 0 + 0 \cdot (-1) = 0
   \]
Calculamos el producto punto de \( \mathbf{v}_4 \) con \( \mathbf{u}_2 \):
\[
   \mathbf{v}_4 \cdot \mathbf{u}_2 = 0 \cdot 1 + 0 \cdot 0 + 1 \cdot 3 + 0 \cdot 2 = 3
   \]
Calculamos el producto punto de \( \mathbf{v}_4 \) con \( \mathbf{u}_3 \):
\[
   \mathbf{v}_4 \cdot \mathbf{u}_3 = 0 \cdot \left(-\frac{1}{3}\right) + 0 \cdot \frac{5}{6} + 1 \cdot 0 + 0 \cdot \frac{1}{6} = 0
   \]
Ahora, calculamos \( \mathbf{u}_4 \) usando estos resultados:
\[
\mathbf{u}_4 = \mathbf{v}_4 - \frac{\mathbf{v}_4 \cdot \mathbf{u}_2}{\mathbf{u}_2 \cdot \mathbf{u}_2} \mathbf{u}_2
\]
Calculamos \( \mathbf{u}_2 \cdot \mathbf{u}_2 \):
\[
\mathbf{u}_2 \cdot \mathbf{u}_2 = 1^2 + 0^2 + 3^2 + 2^2 = 1 + 0 + 9 + 4 = 14
\]
Por lo tanto:
\[
\mathbf{u}_4 = \left(0, 0, 1, 0\right) - \frac{3}{14} \left(1, 0, 3, 2\right)
\]
\[
= \left(0, 0, 1, 0\right) - \left(\frac{3}{14}, 0, \frac{9}{14}, \frac{6}{14}\right)
\]
\[
= \left(0 - \frac{3}{14}, 0, 1 - \frac{9}{14}, 0 - \frac{6}{14}\right)
\]
\[
= \left(-\frac{3}{14}, 0, \frac{5}{14}, -\frac{3}{7}\right)
\]
Con esto, la base ortogonal para $\mathbb{R^4}$ es:
\[
\{\mathbf{u}_1, \mathbf{u}_2, \mathbf{u}_3, \mathbf{u}_4\}
\]
donde:
\[
\mathbf{u}_1 = (2, 1, 0, -1),
\]
\[
\mathbf{u}_2 = (1, 0, 3, 2),
\]
\[
\mathbf{u}_3 = \left(-\frac{1}{3}, \frac{5}{6}, 0, \frac{1}{6}\right),
\]
\[
\mathbf{u}_4 = \left(-\frac{3}{14}, 0, \frac{5}{14}, -\frac{3}{7}\right).
\]
\end{solution}